% \paragraph{Responsabilit�/participation � des projets de recherche nationaux et internationaux depuis 2005.}
% \paragraph{Responsabilit�/participation � des projets de recherche nationaux et internationaux depuis 2012.}
\paragraph{Responsabilit�/participation � des projets de recherche nationaux et internationaux depuis 2015.}
\begin{itemize}
% \item \SR{}{ANR MOSAIC}
%   \item ANR AgriArray~: {\sl New developments towards a genetic and diagnosis use of plant microarrays}, 2006-09;
%   \item ANR AgroBI~: Int�gration d'un centre d'interactions prot�iques (hub) dans le r�seau g�nique global par l'analyse transcriptomique, la mod�lisation des donn�es, et la validation fonctionnelle;
%   \item ANR CoCoGen~: {\sl Comparison of Complete Genomes : an algorithmic and statistical approach to investigate the mechanisms of bacterial genome evolution}, 2008-11;
%   \item ANR NeMo~: {\sl Network Motif}, 2008-10, porteur du projet
%   \item ANR MetaSoil~: {\sl Metagenomic discovery and exploitation of the soil microbial community}, 2008-10;
%   \item ANR CBME~: {\sl Computational Biology for Metagenomics Experiments},
%   2008-10;
  \item PIA Amaizing~: D�velopper de nouvelles vari�t�s de ma�s pour une
  agriculture durable, 2010-17;
%   \item ANR CNV-Maize~: Etude d'association sur g�nome entier entre variation structurale, variation des caract�res d'int�r�t agronomique et h�t�rosis chez le ma�s, 2011-14;
%  \item PIA ABS4NGS~: {\sl Algorithmic, Bioinformatics and Statistics for Next Generation Sequencing}, 2011-15, responsable du WP {\sl transcriptional landscape}
%  \item INCa\footnote{Institut national de recherche sur le cancer} EPIREG~: Identification des r�gions chromosomiques inactiv�es ou activ�es par des m�canismes �pig�n�tiques, 2011-2014
%  \item ONEMA\footnote{Office national de l'eau et des milieux aquatiques}~: D�veloppement d'outils de bio-indication "phytobenthos" et "macro-invert�br�s benthiques" pour les cours d'eau de Mayotte, 2013-15
 \item ANR HydroGene~: Comparative Metagenomic for Measuring Biodiversity,  
Application to Ocean Life Studies, 2014-2018
 \item ANR NGB~: Next generation biomonitoring, 2016-2020 
 \item ANR EcoNet~: Advanced statistical modelling of ecological networks, 2017-2021 
\end{itemize}
J'ai �galement particip� � des projet soutenus plus sp�cifiquement par l'INRA
\begin{itemize}
\item Learn-Biocontrol~: Learning microbial networks from metabarcoding data; application to biological control (soutenu par le m�ta-programme MEM), 2016-2017
\item CNV4sel~: Strategies for CNV discovery in animal and plant species and their use in
breeding (soutenue par le m�ta-programme SelGen), 2016-2018
% \item BrassicaDiv-Patho~: Analyse de la diversit� et des r�seaux d�associations des communaut�s bact�riennes, fongiques et virales associ�es � {\sl brassica napus} dans un contexte pathobiome soutenu par le m�ta-programme MEM), 2017-2018
\item BCMicrobiome~: Learning microbial networks from metabarcoding data: application to biological control (soutenu par le consortium Biocontr�le), 2017-2020
\item KineTicks~: Network and modelling analyses to describe the dynamics of {\sl Ixodes ricinus} microbiome and its influence in pathogen dynamics (soutenu par le m�ta-programme MEM), 2018-2020
\end{itemize}

% \paragraph{Projet en cours d'�valuation.}
% \begin{itemize}
%  \item H2020 ENABLORE~: {\sl Efficient Novel Analytic and Bioinformatic tools for Large-scale Omics in Research and Enterprise}, responsable du WP {\sl Refined (ultra) high dimensional modelling, omics integration}.
% \end{itemize}
